\documentclass[10pt,a4paper]{amsart}
\usepackage[margin=19mm]{geometry}
\usepackage{amsmath,amssymb,mathtools}
\usepackage[scr=rsfs,cal=euler]{mathalfa}
\usepackage{xfrac}
\usepackage[unicode,hidelinks,bookmarks=false]{hyperref}
\newcommand{\ie}{i.e.\ }
\newcommand{\cf}{cf.\ }
\newcommand{\resp}{respectively\ }
\newcommand{\Subsection}{\subsection}
\providecommand{\nobreakdash}{\nobreak\hbox{-}}
\setlength{\emergencystretch}{2em}
\multlinegap0pt
\usepackage{amssymb}
\usepackage{dsfont}
\usepackage{mathtools}
\usepackage[shortlabels]{enumitem}
\usepackage{tikz}
\newtheorem{theorem}{Theorem}
\newtheorem{lemma}{Lemma}
\title{Word-Representation of Melon Graphs}
\author{Khyodeno Mozhui, K. V. Krishna}
\date{Source: arXiv:2602.14179v1 (CC BY 4.0)}
\begin{document}
\maketitle

\noindent\emph{Source and licence.} Word-Representation of Melon Graphs. Version arXiv:2602.14179v1 (CC BY 4.0), distributed by arXiv under the Creative Commons Attribution 4.0 International licence, http://creativecommons.org/licenses/by/4.0/, as recorded in the arXiv licence metadata for this exact version and shown on the abstract page https://arxiv.org/abs/2602.14179v1. Full licence notice: \texttt{licenses/arxiv-2602.14179-CC-BY-4.0.txt}.\par\smallskip

\noindent\textit{Editorial scope.} Counted unit: the article's theorem 00'adj\_fc (source0.tex lines 521--523). The article's complete original proof of it is delivered with it (source0.tex lines 525--567, one \texttt{proof} environment of 43 lines). No mathematics is written, repaired or re-derived: the statement and the proof are the article's own lines, in the article's own notation, and no retained line is substituted or re-flowed.  Retained uncounted as the source-local context the proof needs: lines 341--347; lines 358--362; lines 392--394.  Nothing was omitted from any retained span, so the package carries no omission record.  No retained line contains a citation, so the wrapper carries no bibliography and no bibliography entry.  No retained line contains a figure environment, an image-inclusion command or drawing code, so no image is delivered and the package declares no asset dependency.  Editorial formatting only, in addition: an AMS \texttt{amsart} preamble that restates the article's own declarations and macros used by the retained lines, namely the environments \texttt{theorem}, \texttt{lemma} on separate counters, together with the standard packages those lines need.  The retained environments print in this document as Theorem 1, Lemma 1 in order of appearance, whereas the article prints the same items with the numbering of its own full document.  The complete native source of the article as distributed in the arXiv e-print for this exact version is bundled unchanged as \texttt{sources/194694/source0.tex}.
\begin{theorem}\label{chr_com_file}
	A melon graph $M = (E_1, \ldots, E_m)$ is a comparability graph if and only if one of the following holds: 
	\begin{enumerate}
		\item All constituent paths of $M$ have the same parity.
		\item The endpoints $0$ and $0'$ are adjacent, and all constituent paths of even parity are of length two. 
	\end{enumerate}	
\end{theorem}

For $k \ge 2$, let $u$ and $v$ be the words over the vertices $c_2, c_3, \cdots, c_{2k-1}$ given by 
\begin{eqnarray}
	\label{eq-1} u &= c_{4} c_{3} \cdots c_{2i} c_{2i-1} \cdots c_{2k-2} c_{2k-3},\; (\text{for clarity, } 2 \le i \le k-1);\\
	\label{eq-2} v &= c_{2k-2} c_{2k-1} \cdots c_{2j} c_{2j+1} \cdots c_2 c_3,\; (\text{for clarity, } 1 \le j \le k-1).
\end{eqnarray} 

\begin{lemma}\label{even_cycle_word}
	For $k \ge 1$, let $\langle 0', c_1, c_2, \ldots, c_{2k}, 0 \rangle$ be a path on $2k+2$ vertices. If $0$ and $0'$ are adjacent, then the corresponding even cycle $C$ is  represented by the concatenation of the permutations $w_1 = 0' c_2 c_1 uc_{2k} c_{2k-1} 0$, $w_2 = c_{2k} v 0' c_1 0$ and $w_3 = 0' c_{2k} 0 v c_1$, where $u$ and $v$ are the words given in equations (\ref{eq-1}) and (\ref{eq-2}), respectively.
\end{lemma}

\begin{theorem}\label{00'adj_fc}
	Let a melon graph $M_c$ be a comparability graph. If the endpoints $0$ and $0'$ are adjacent in $M_c$, then $\mathcal{R}^p(M_c) \le 3$.
\end{theorem}

\begin{proof}
	Without loss of generality, let $M_c = (E_1,E_2, \ldots, E_m)$ such that $|E_1| \ge |E_2|\ge \cdots \ge |E_m|$, where $|E_i|$ is the number of vertices of $E_i$.  Since $E_m$ is of length one, note that the endpoints $0$ and $0'$ are adjacent in $M_c$. For $i \ge 2$, if a constituent path $E_i$ is of odd length, then the induced subgraph $M_c[E_i]$ is an even cycle. If $E_i$ is of even length, then it should be of length two (by Theorem \ref{chr_com_file}) so that $M_c[E_i]$ is a triangle. In the construction of permutations on the vertices of $M_c$, we require the intermediate vertices, say $x_i$ and $y_i$, adjacent to $a_i$ and $a_i'$ in $E_i$, respectively. 
	
	\noindent\textbf{Part-I.} Suppose all the constituent paths of $M_c$ are of odd parity.
	For $1 \le i < m$, note that $E_i$ is of length at least three. Consider the permutations $w_1 = 0'y_ia_i' u_i a_i x_i 0$, $w_2 = a_i v_i 0' a_i' 0$ and $w_3 = 0' a_i 0 v_i a_i'$, where $u_i$ and $v_i$ are permutations on the remaining intermediate vertices of $E_i$, as stated in Lemma \ref{even_cycle_word}. 
	Construct the following three permutations on the vertices of $M_c$: 
	\begin{align*}
		p_1' &= 0' y_{m-1} a_{m-1}' \cdots y_i a_i' \cdots y_1 a_1' u_{1} \cdots u_i  \cdots  u_{m-1}    a_1 x_1 \cdots a_i x_i   \cdots  a_{m-1} x_{m-1} 0 \\
		p_2' &= a_{m-1}v_{m-1}        \cdots  a_iv_i \cdots a_1 v_1    0'  a_{m-1}'    \cdots a_{i}' \cdots a_1'0\\
		p_3' &= 0'a_{m-1} \cdots a_{i} \cdots a_1 0 v_{1}a_{1}' \cdots v_ia_i' \cdots      v_{m-1} a_{m-1}',
	\end{align*} (for clarity, $1 \le i \le m-1$).
	If $E_i$ if of length three, note that $x_i = a_i'$ and $y_i = a_i$. Then, put $v_i = \epsilon$ in the permutations $p_2'$ and $p_3'$, and put $a_i' = y_i = \epsilon$ in $p_1'$. 
	
	Using induction on $m$, we show that two vertices are adjacent in $M_c$ if and only if they alternate in $p_1'p_2'p_3'$. 
	The statement holds for $m = 1$ by Lemma \ref{even_cycle_word}. Assume that the statement holds if $M_c$ has at most $m - 1$ constituent paths. It is sufficient to show that, for $a \in  E_2 \cup \ldots \cup E_{m-1} \cup E_m$ and $b \in E_{1} \setminus \{0,0'\}$, $a$ and $b$ are adjacent in $M_c$ if and only if $a$ and $b$ alternate in $p_1'p_2'p_3'$.
	
	Suppose  $a$ and $b$ are adjacent in $M_c$. Then, clearly $a \in \{0',0\}$ and $b \in \{a_{1}', a_{1}\}$. If $a = 0$, then $b = a_{1}$; in this case, we have $ba \ll p_1'$, $ba \ll p_2'$, and $ba \ll  p_3'$. In case $a = 0'$, then $b = a_{1}'$ and $ab \ll p_1'$, $ab \ll p_2'$, and $ab \ll p_3'$. Hence, in any case, $a$ and $b$ alternate in $p_1'p_2'p_3'$.
	
	
	Conversely, suppose that $a$ and $b$ are not adjacent in $M_c$. Then, $a$ must be an intermediate vertex of some path $E_i$ ($2 \le i \le m-1$). We show that  $a$ and $b$ do not alternate in $p_1'p_2'p_3'$ in the following three cases:
	
	\begin{itemize}
		\item $b \in \{a_1', y_1\}$. If $a \in \{a_i', y_i\}$, then $ab \ll p_1'$ and $ba \ll p_3'$. If $a \not  \in \{a_i', y_i\}$, then $ba \ll p_1'$ and $ab \ll p_2'$.
		
		\item $b \in \{a_1,x_1\}$. If $a  \in \{a_i, x_i\}$, then $ba \ll p_1'$ and $ab \ll p_2'$. Else if $a  \not \in \{a_i, x_i\}$, $ab \ll p_1'$ and $ba \ll p_3'$.
		
		\item $b \not \in \{a_1,a_1',x_1,y_1\}$, i.e., $b \ll u_1$ and $b \ll v_1$. If $a  \in \{a_i', y_i\}$, then $ab \ll p_1'$ and $ba \ll p_3'$. Else if $a \not \in \{a_i',y_i'\}$, then $ba \ll p_1'$and $ab \ll p_2'$.
	\end{itemize}
	Hence, if all the constituent paths of $M_c$ are of odd parity,  then $M_c$ can be represented by the word $p_1'p_2'p_3'$. 
	
	\noindent\textbf{Part-II.} Suppose that the melon graph $M_c$ has constituent paths of even length. Let $t$ be the least index such that $E_t$ is of length two. Then, the constituent paths  $E_t, E_{t+1} \ldots, E_{m-1}$ are of length two, with intermediate vertices, say $a_t, a_{t+1}, \ldots, a_{m-1}$, respectively. Construct the following three permutations on the vertices of $M_c$:
	\begin{align*}
		q_1' &= 0' y_{t-1} a_{t-1}' \cdots y_j a_j' \cdots y_1 a_1'  u_{1} \cdots u_j  \cdots  u_{t-1}    a_1 x_1 \cdots a_j x_j   \cdots  a_{t-1} x_{t-1}  a_{m-1} \cdots a_{i} \cdots a_{t} 0 \\
		q_2' &=  a_{t-1}v_{t-1}        \cdots  a_jv_j \cdots a_1 v_1    0'a_{t-1}'     \cdots a_{j}' \cdots a_1' a_t \cdots a_{i} \cdots a_{m-1} 0\\
		q_3' &= 0'  a_t \cdots a_{i} \cdots a_{m-1} a_{t-1} \cdots a_{j} \cdots a_1 0 v_{1}a_{1}' \cdots v_ja_j' \cdots      v_{t-1} a_{t-1}'.
	\end{align*} (for clarity, $t \le i \le m-1$ and $1 \le j \le t-1$). 
	For $F = E_1 \cup E_2 \cup \cdots \cup E_{t-1} \cup E_{m}$, observe that $(q_1'q_2'q_3')|_F = (p_1'p_2'p_3')|_F$. Thus, by Part-I, $(q_1'q_2'q_3')|_F$  represents $M_c[F]$ permutationally. It remains to show that, for $a = a_i$ ($t \le i\le m-1$) and $b \in E_j$ ($1 \le j \le m$), $a$ and $b$ are adjacent if and only if $a$ and $b$ alternate in $q_1'q_2'q_3'$.
	
	Suppose $a$ and $b$ are adjacent in $M_c$. Then, $b \in \{0,0'\}$. If $b = 0$, then clearly $ababab \ll q_1'q_2'q_3'$. Otherwise, $bababa \ll q_1'q_2'q_3'$. Thus, $a$ and $b$ alternate in $q_1'q_2'q_3'$. Conversely, suppose $a$ and $b$ are non-adjacent in $M_c$. If $a = a_i$ and $b = a_j$ (for $t \le i < j \le m-1$), then $a_j a_i \ll q_1'$ and $a_ia_j \ll q_2'$. If $a = a_i$ (for $t \le i \le m-1$) and $b \in E_1 \cup E_2 \cdots \cup  E_{j-1}$, then $ba\ll q_2'$ and $ab \ll q_3'$. Thus, $a$ and $b$ do not alternate in $q_1'q_2'q_3'$.  
	Hence, if $M_c$ has constituent paths of even length, then $M_c$ can be represented by the word $q_1'q_2'q_3'$.
	
	Hence, if the endpoints $0$ and $0'$ are adjacent in $M_c$, then $\mathcal{R}^p(M_c) \le 3$.	\qed
\end{proof}

\end{document}
